\section{Diagrama de arquitectura del sistema}

El sistema híbrido propuesto se compone de tres capas principales:

\begin{enumerate}
    \item \textbf{Capa de entrada (Frontend):} Interfaz donde el usuario ingresa la URL o el texto de la noticia financiera a verificar.
    \item \textbf{Capa de clasificación (LLM):} Un modelo de lenguaje de gran escala clasifica la noticia en tres categorías (verdadera, dudosa, falsa) y emite una puntuación de confianza.
    \item \textbf{Capa de verificación (Agente):} Cuando la confianza del clasificador es inferior al umbral definido, un agente inteligente basado en ReAct consulta fuentes institucionales colombianas y emite un veredicto fundamentado.
\end{enumerate}

El flujo completo es orquestado por LangGraph mediante un grafo de estados con nodos condicionales.

\section{Listado de fuentes institucionales del agente}

\begin{itemize}
    \item Superintendencia Financiera de Colombia (\url{https://www.superfinanciera.gov.co})
    \item Banco de la República (\url{https://www.banrep.gov.co})
    \item Departamento Administrativo Nacional de Estadística -- DANE (\url{https://www.dane.gov.co})
    \item ColombiaCheck (\url{https://colombiacheck.com})
    \item Portafolio (\url{https://www.portafolio.co})
    \item La República (\url{https://www.larepublica.co})
\end{itemize}
